\section{Closed-Loop Simulations (SIL)}
\label{sec:sil}

\subsection{Remark on Trim Commands}
\label{sec:sil-trim}
% Nota pedida pela orientadora: deixar clara a questao do delta-comando de
% trimagem vs deltas dados nas simulacoes.
All controllers in this work operate on \emph{incremental} commands about
the trim condition of Section~\ref{sec:trim}: the compensators output
$\Delta\delta$, and the trim commands $\delta^e$
of~\eqref{eq:dtrim} are added before the actuator models. Consequently,
whenever a simulation reports a command signal, the reader should
distinguish the absolute command $\delta = \delta^e + \Delta\delta$ from
the increment $\Delta\delta$ produced by the controller. In steady flight
at the trim point, the absolute commands settle at $\delta^e$ (e.g., the
elevator at $7.62^\circ$), not at zero.

\subsection{Comparison Protocol and Known Asymmetries}
\label{sec:protocol}
Every head-to-head result in this article follows the same protocol
--- in the spirit of the experimental controller benchmarks of the
recent literature~\cite{okasha2022design,rinaldi2023comparative}: the
\emph{same} nonlinear plant, the \emph{same} reference amplitudes
and timings, and the \emph{same} metric windows, with each control law
running in its own design harness, exactly as delivered. That last
clause matters, because the two harnesses are not identical, and an
honest comparison must say so. Three asymmetries are known and
uncontrolled:
\begin{enumerate}
  \item \emph{Actuator models.} The PID harness includes the full
        actuator chain of Section~\ref{sec:model} (position and rate
        saturations, first-order servo and engine lags). The LQRy
        harness includes first-order actuator lags --- 24\,rad/s on
        the surfaces versus the 20\,rad/s assumed on the PID side,
        and 0.1\,s versus 0.30\,s for the throttle --- and a $[0,1]$
        throttle saturation, but omits the rate limits and surface
        position saturations of the PID harness.
  \item \emph{Reference shaping.} The PID applies the
        $1/(3s+1)$ prefilter of Section~\ref{sec:pid} to every
        reference increment; the LQRy receives raw reference steps.
  \item \emph{Protection logic.} The $\pm10^\circ$ $\theta^{ref}$
        clamp and its anti-windup exist only in the PID cascade.
\end{enumerate}

To bound the effect of these asymmetries on the conclusions, the PID
scenarios were repeated in an \emph{equalized} configuration ---
actuator bandwidth matched to 24\,rad/s, rate limits removed, clamp
and prefilter disabled, matching the LQRy harness.
Fig.~\ref{fig:eq-comb} overlays the two PID configurations on the
combined scenario: the equalization shifts the guidance captures
moderately --- airspeed rise time from 9.5 to 5.8\,s and settling from
18.0 to 11.4\,s, heading settling from 33.0 to 29.3\,s, altitude
settling from 44.4 to 40.8\,s, all with zero overshoot --- without
changing any comparative conclusion drawn below.

\begin{figure}[!t]
\centering
\includegraphics[width=\linewidth]{fig_eq_comb.png}
\caption{Equalization ablation, combined scenario: PID as delivered
vs.\ PID with the harness equalized to the LQRy rig (24\,rad/s
actuators, no rate limits, no clamp, no prefilter), with the LQRy
response for context.}
\label{fig:eq-comb}
\end{figure}

Two further implementation notes complete the disclosure. First, an
earlier fixed-gain LQRy build fed back the \emph{absolute} state
instead of the deviation from trim, producing a start-up kick (initial
elevator command of $4.11^\circ$ instead of the trim $7.62^\circ$);
the gain-scheduled build compared here feeds back trim deviations by
construction, and its runs were verified to start at the trim command.
The embedded (incremental-form) implementation of
Section~\ref{sec:hil} was immune to this distinction throughout.
Second, the LQRy rig initializes the plant at $\VT=12.0$\,m/s while
its stored controller trim corresponds to 11.62\,m/s, so LQRy runs
begin with a small longitudinal offset that the airspeed loop resolves
before the first reference event.

\subsection{PID on the Nonlinear Model}
\label{sec:sil-pid}
The cascade PID was evaluated on the nonlinear model under three
maneuvers, with the linear closed-loop response overlaid in each case:
a longitudinal climb ($h$ $+10$\,m), a lateral--directional heading change
($\psi$ $+15^\circ$), and an aggressive combined maneuver ($h$ $+30$\,m and
$\psi$ $+40^\circ$ commanded simultaneously, saturating the
$\theta^{ref}$ clamp for ${\approx}10$\,s).
Fig.~\ref{fig:nl-lin} shows the combined case. Linear and nonlinear
responses agree closely even in the aggressive maneuver (maximum
deviations of 0.38\,m in altitude, $0.52^\circ$ in pitch, and
$0.27^\circ$ in bank; below 0.21\,m/$0.75^\circ$/$0.05^\circ$ in the
single-axis maneuvers), validating both the linearization and the
linear-design--nonlinear-verification workflow. This validation was
performed with the original loop-shaping tuning; its conclusion ---
the fidelity of the linearization --- is independent of the
subsequent requirements-driven retune.

\begin{figure}[!t]
\centering
\includegraphics[width=\linewidth]{fig_nl_lin_agressiva.png}
\caption{PID, aggressive combined maneuver ($h$ $+30$\,m and $\psi$
$+40^\circ$ simultaneous): nonlinear model vs.\ linear closed loop.}
\label{fig:nl-lin}
\end{figure}

\subsection{LQRy on the Nonlinear Model}
\label{sec:sil-lqry}
The LQRy was evaluated directly on the nonlinear model, in its own
harness, under the same maneuver set used for the comparison below.
Its closed-loop character is consistent across maneuvers: fast,
high-bandwidth captures --- airspeed and altitude in ${\approx}6$\,s
with essentially no overshoot, heading in ${\approx}13$\,s with
visible overshoot (Table~\ref{tab:metrics}). Unlike the PID case, no linear-overlay
validation is presented: the linear design models of
Section~\ref{sec:lqry} are the basis of the synthesis, and the
nonlinear evaluation here is the verification step of record for the
delivered design.

\subsection{Controller Comparison}
\label{sec:sil-comparison}
The evaluation maneuvers deserve a word of justification. Doublets are
the conventional, well-characterized excitation of flight-test
practice --- originally for identification, where their frequency
content is matched to the modes of
interest~\cite{morelli1999flight,klein2006aircraft} --- and step and
doublet captures underlie the demonstration requirements of the
flying-qualities specifications~\cite{milf8785c1980}. For
\emph{unmanned} aircraft, however, the literature argues that
mission-based evaluation is the most transferable criterion, since
manned handling-qualities metrics do not carry over
directly~\cite{greene2014toward}; command-profile tracking of
airspeed, altitude, and attitude remains the standard validation
practice for small-UAV control laws~\cite{shen2025model}. The
campaign therefore uses both: per-axis guidance scenarios to
characterize each channel, and a 300\,s \emph{combined mission}
(airspeed step, heading doublet, altitude doublet, overlapping in
time) as the mission-level verdict.

The two controllers are compared on the nonlinear model under four
guidance-level scenarios
(\emph{VelHold}+\emph{PsiHold}+\emph{AltitudeHold}) with
\emph{identical} reference profiles on both rigs (amplitudes injected
into each rig's unmodified reference harness):
\begin{enumerate}
  \item \emph{Airspeed}: step $12\!\to\!15.2$\,m/s at 5\,s;
  \item \emph{Heading}: doublet of $\pm5^\circ$ at 50/100/150\,s;
  \item \emph{Altitude}: doublet of $\pm5$\,m at 80/160/240\,s;
  \item \emph{Combined mission}: all three simultaneously, 300\,s
        (flight case~4 of the campaign of Section~\ref{sec:hil}).
\end{enumerate}
Table~\ref{tab:metrics} reports the classical step-response metrics
per guidance channel, and Fig.~\ref{fig:cen-comb} shows the
combined mission.
\begin{figure}[!t]\centering
\includegraphics[width=\linewidth]{fig_alt_pid_lqry.png}
\caption{Altitude channel of the combined mission (SIL): reference,
PID, and LQRy on a single axis.}\label{fig:alt-single}\end{figure}
The roles observed in the earlier fixed-gain iteration of this
comparison are reversed: the gain-scheduled LQRy is now the fast,
high-bandwidth law --- it captures heading in ${\approx}13$\,s at the
cost of $16.4\%$ overshoot, and completes the airspeed and altitude
captures in ${\approx}6$\,s with essentially no overshoot --- while
the requirements-retuned PID is the smooth one, with zero overshoot on
every channel and settling times of 18--44\,s. Neither tuning
dominates: the LQRy buys speed at the cost of a heading-overshoot
requirement violation; the PID buys strict compliance at the cost of
bandwidth.

\begin{table}[!t]
\caption{Per-axis guidance captures (SIL, identical references),
metrics on the $+$step segment of each channel. OS = overshoot, $t_r$
= 10--90\% rise time, $t_s$ = 2\% settling time, $e_{ss}$ =
steady-state error.}
\label{tab:metrics}
\centering
\resizebox{\linewidth}{!}{%
\begin{tabular}{lcccccccc}
\toprule
 & \multicolumn{4}{c}{PID} & \multicolumn{4}{c}{LQRy} \\
Ch. & OS[\%] & $t_r$[s] & $t_s$[s] & $e_{ss}$[\%] & OS[\%] & $t_r$[s] & $t_s$[s] & $e_{ss}$[\%] \\
\midrule
$\VT$ & 0.0 & 9.48 & 18.0 & 0.00 & 0.0 & 3.04 & 5.8 & 0.00 \\
$\psi$ & 0.0 & 16.14 & 33.0 & 0.50 & 16.4 & 0.83 & 12.6 & 0.01 \\
$h$ & 0.0 & 23.37 & 44.4 & 0.06 & 0.6 & 3.44 & 5.9 & 0.19 \\
\bottomrule
\end{tabular}}
\end{table}

\begin{figure}[!t]
\centering
\includegraphics[width=\linewidth]{fig_cen_comb.png}
\caption{Scenario 4 (combined mission): PID vs.\ LQRy on the nonlinear
model, identical references.}
\label{fig:cen-comb}
\end{figure}

Table~\ref{tab:compliance} closes the loop with the requirements of
Section~\ref{sec:requirements}, evaluated on the guidance channels of
the combined mission. The PID meets every row. The LQRy settles far
inside the 60\,s limit on every channel (5.8--7.0\,s) and meets the
steady-state row, but violates the overshoot requirement on the
heading channel ($24.5\%$ against the $10\%$ line).

\begin{table}[!t]
\caption{Requirements compliance on the combined mission (SIL),
guidance channels, identical references.}
\label{tab:compliance}
\centering
\begin{tabular}{lcccccc}
\toprule
 & \multicolumn{3}{c}{PID} & \multicolumn{3}{c}{LQRy} \\
Requirement & $\VT$ & $\psi$ & $h$ & $\VT$ & $\psi$ & $h$ \\
\midrule
OS $<10\%$ & 0.0 & 0.0 & 0.0 & 0.2 & \textbf{24.5} & 0.2 \\
$t_s<60$\,s & 18.0 & 38.1 & 45.1 & 5.8 & 6.9 & 7.0 \\
$e_{ss}\le5\%$ & 0.00 & 0.72 & 0.07 & 0.00 & 0.00 & 1.75 \\
\bottomrule
\end{tabular}
\end{table}

The embedded (HIL) campaign of Section~\ref{sec:hil} was executed with
the \emph{earlier fixed-gain} LQRy build (Section~\ref{sec:lqry}), so
its LQRy-side numbers characterize that deployment rather than the
gain-scheduled law compared here; the PID-side SIL--HIL fidelity
results are unaffected.

% (Robustez movida para secao propria: sections/05_robustness.tex.
%  Remark de execution order movido, encurtado, para o fim de
%  sec:determinism no HIL — resolve o "???" da orientadora.)
